\section{Exploratory note on prior strips and an exponential tilt}
This note records a route to a further numerical improvement. It is
separate from the proved cosine bound. Its independent computation
checks the last integral conditionally on the analytic comparison below.

\paragraph{The normalization.}
Write $L=\log|T|$, $s=1+u/L+i\ell T$, with $u>0$ fixed, and consider a
real compactly supported autocorrelation $h$ as in the preceding section.
Suppose that the standard factors at positive frequency having positive
weight in the amplifier have no zero with
$(1-\beta)L<a$ at distance at most $1$ from $\ell T$, uniformly for
all sufficiently large $T$. Choose the inequality with a fixed margin.
For $W>0$, set
$q(x)=e^{(u+a)x}h(Wx)$ for $x\ge0$ and extend it evenly.
Then $q\ge0$ and $q(0)=1$. At real abscissa $1+u/L$ the normalized
weighted Euler sum is finite and all prime coefficients still have the
original signs. The kernel for a zero $\rho$ is
\[
 K_q(s-\rho)= L\int_0^\infty q(x)e^{-L(s-\rho)x}\,dx .
\]
Its real part is nonnegative for the zeros in the indicated disk:
the effective real argument of the Laplace transform of $h(Wx)$ is
$L(1-\beta)-a\ge0$. This follows from positive definiteness after
multiplication by $e^{-v|x|}$ with $v\ge0$. The main target contributes
$L H((w-a)/W)/W$ when $w=L(1-\beta)$; the static pole contributes
$L H(-a/W)/W$.

\paragraph{A comparison with the ordinary disk estimate.}
The Fourier density $d\nu(v)=\widehat q(v)\,dv/(2\pi)$ gives a signed
finite measure of total mass one. For the cosine profile, $q$ is
piecewise twice continuously differentiable with bounded first derivative
and integrable second derivative on the positive half-axis. The even
extension can have a corner at zero. Two integrations by parts show
$|d\nu(v)|\ll(1+|v|)^{-2}\,dv$. For $\Re z>0$ they also give the
uniform estimate
\[
              |K_q(z)-z^{-1}|\le
       \frac{|q'(0+)|+\|q''\|_{L^1(0,\infty)}}{L|z|^2}.
\]
An approach to retaining the Weyl cost is to average the full Hadamard
identity over $s+iv/L$ using $\nu$, and subtract the unweighted
Hadamard identity. In the difference, every zero replaces the simple
pole by $K_q(z)-1/z$. For any fixed $\eta>0$, all zeros with
$|s-\rho|\ge\eta$ then cost $O_{\eta}(1)$: compare
$|s-\rho|^{-2}$ to $|2+i\ell T-\rho|^{-2}$ and use the positive
Hadamard identity at real part $2$, which bounds the latter sum by
$O(L)$. For each fixed $T$ absolute Fubini against $|\nu|$ follows
from the same positive identity at $1+u/L$, Euler absolute convergence
and $\int \log(2+|v|)\,d|\nu|(v)<\infty$.
Each shifted gamma term changes by $o(L)$ since $\nu$ has mass one.
For $|v|\le L|T|/2$ and $\ell\ge1$, the two digamma real parts
are $\tfrac12\log|T|+O(1)$. On the complement the variation contributes
$O((\log(2+L|T|)+L)/(L|T|))$ using the displayed density decay and
the elementary uniform logarithmic digamma bound. Fixed static terms
have bounded gamma average. Poles are kept as exact kernels.

In the original, unaveraged disk Poisson--Jensen upper bound, retain
only zeros with $|s-\rho|<\eta<1/4$. The terms for the other
in-disk zeros have a favorable sign before the difference is added.
The radius-$R$ correction for the remaining zeros is
$O(\eta L)$ for $R\in[1/4,1/3]$, because the standard local count is
$O(L)$. After combining with the signed difference, the test kernel
is nonnegative at every retained zero. Let $T\to\infty$ and then
$\eta\downarrow0$. This comparison leads to the proposed inequality
\[
           H((w-a)/W)-P H(-a/W)\le WG .
\]
The order of operations is important: signed averaging of the disk
upper bound itself does not preserve that bound.

\paragraph{Scalar and static terms.}
For signed static coefficients, each fixed completed factor is
nonvanishing in a fixed punctured neighborhood of $1$. Remove any
actual pole first. The integrated regular contribution is $O(1)$ by
the far-zero difference estimate and the original bounded logarithmic
derivative. At high scalar frequencies, including any negative-weight
terms, an absolute bound $o(L)$ for the real part of the weighted scalar logarithmic
derivative is needed. For the Riemann zeta function the usual
Vinogradov--Korobov width and log-derivative bound provide this:
$|\Re(\zeta'/\zeta)(1+u/L+it)|=o(L)$ uniformly when
$|t|\asymp |T|$, by the classical bounds in the zero-free region.
On $|v|>L|T|/2$ absolute convergence at $1+u/L$ costs $O_u(L)$,
and the variation there is $O((L|T|)^{-1})$. The high scalar poles
are part of this direct bound. For fixed Hecke characters use the real-part estimate in the
fixed-rank argument, from Coleman, \emph{A zero-free region for the
Hecke $L$-functions}, Theorem~2, and the positive Hadamard comparison
there. The estimate is uniform for $|t|\asymp|T|$ because $\log|t|=L+O(1)$.

\paragraph{A numerical test in the elliptic case.}
For an elliptic symmetric power of index $m\ge1$, the auxiliary
isobaric triple $\{0,m,2m\}$ of the earlier zero-free argument has
three static poles and target multiplicity four. The total degree at
shifted height is $4(m+1)^2$. Ordinary Stirling plus prime positivity
gives $4/(u+w)\le3/u+2(m+1)^2+o(1)$, where the variables are scaled by
$L$. Optimizing over $u$ yields the prior width
$2(1-\sqrt{3/4})^2/(m+1)^2$. For $1\le m\le5$ this exceeds
$0.0009$ by a strict margin. The uniform fixed-rank theorem gives
stronger priors of $1/(6.348n)$ for rank $n$.
The degree-twenty amplifier has largest moving rank six and the
original quartic square has largest moving rank four. Thus $a=0.026$
is available for the former and $a=0.03937$ for the latter.

For any fixed $k,S>0$ with $kS<\pi$, set $c=\cos(kS)$ and take
$g(x)=(\cos(kx)-c)\mathbf1_{[-S,S]}(x)$. For $0\le t\le2S$ its
unnormalized autocorrelation is
\[
 B(t)=(2S-t)\left(c^2+\tfrac12\cos kt\right)
       +\frac{\sin(k(2S-t))}{2k}
       -\frac{2c}{k}\{\sin(kS)+\sin(k(S-t))\}.
\]
This follows by integrating the product of the two cosines on
$[-S,S-t]$. Normalize by $B(0)$ and compute $H$ over $[0,2S]$.
For degree twenty take $k=2.35698$, $S=0.78646$, $a=0.026$ and
$W=1/6.005086$. The independently enclosed scalar surplus
$H(1-a/W)-PH(-a/W)-WG$ exceeds $2.126\cdot10^{-8}$.
For the quartic square take $k=2.28308$, $S=0.78777$, $a=0.03937$
and $W=1/5.943825$. Its corresponding surplus exceeds
$2.635\cdot10^{-8}$. The check uses rational parameters and interval
quadrature of entire functions on $[0,1]$.
See \texttt{tilted\_check.py}. These scalar computations are separate
from the analytic comparison above. That comparison has a human
argument and has not been translated to Lean here.
